\subsection{Model Fit}

Our model estimates generate predictions that match our data well along multiple dimensions. To demonstrate this, Table \ref{tab:fit_c} presents a comparison of key moments in our data and their predicted model analogues for claims occurrence, monitoring scores and pricing factors. Table \ref{tab:fit_d} presents a comparison of moments for coverage choice shares, monitoring opt-in share for new customers, and attrition shares for renewal customers. This includes moments that capture selection patterns, i.e., how claim costs vary across alternatives, for both coverages and monitoring opt-in.\footnote{We do not observe claims outcomes for attrited consumers.} In each case, our model predictions are able to recover the distributions of outcomes observed in our data to a high level of granularity.
Furthermore, the model replicates choice shares and selection patterns across major product-menu changes, including the introduction of the monitoring program and a change in the mandatory minimum coverage requirement in Illinois from \$40,000 to \$50,000.

\subsection{Risk Rating and Selection}

Our estimates allow us to directly quantify and visualize the effectiveness of the firm's risk rating with and without monitoring, as well as consumers' selection into the program.

\clearpage
\enlargethispage{3em}
\begingroup
\let\oldtable\table
\def\table[#1]{\oldtable[H]}
\inputifexists{\exhibitpath/tab_4}
\vspace{-2em}
\inputifexists{\exhibitpath/tab_5}
\endgroup
\clearpage

\paragraph{Risk Rating}

The firm's risk rating capability is captured by our pricing and monitoring score models. The former captures how prices evolve dynamically (from the observed initial level) based on consumers' risk class and, if available, their monitoring scores. The latter directly connects the scores with consumers' claim risk.

Figure \ref{fig:risk_rating_selection}a shows a strong positive correlation between the first renewal price that a consumer faces and their log claim risk ($\lambda$). Monitored consumers are subject to a steeper pricing curve, in which consumers with lower risk types see lower prices and vice versa. The resulting advantageous selection can be seen in Figure \ref{fig:risk_rating_selection}b. Monitored drivers face a stochastically lower claim risk distribution (in the first-order sense) than unmonitored drivers. Our data and estimation cover several baseline pricing regimes and monitoring regimes. The figures above focus on the latest pricing and monitoring regime, which serves as the baseline for our counterfactual analyses.  \Cref{fig:R-vs-risk-by-regime} compares risk rating across regimes.

\begin{figure}[htbp!]
\begin{centering}
\begin{tabular}{cc}
\includeifexists[scale=0.6]{\exhibitpath/fig_6a.png} &\hspace{0.05cm} \includeifexists[scale=0.5]{\exhibitpath/fig_6b.png}\tabularnewline
(a) &\hspace{0.05cm} (b)
\end{tabular}
\caption{Risk Rating and Advantageous Selection into Monitoring \label{fig:risk_rating_selection}}
\par\end{centering}
\vspace{0.1cm}
\begin{minipage}{\textwidth} % choose width suitably\vspace{0.25cm}
    {\footnotesize \emph{Notes: }(a) is a binned scatter plots first-renewal price (with and without monitoring) against consumers' estimated Poisson claim risk. We focus on the first-renewal price of the standard \$50,000 coverage plan.
    (b) plots the density of the Poisson claim rate by monitoring opt-in.
    \par}
\end{minipage}

\vspace{1.5\baselineskip} % small gap
\end{figure}
\inputifexists{\exhibitpath/tab_6}

\subsection{Parameter Estimates\label{subsec:estimation_results}}

Table \ref{tab:model_params_summary} presents the empirical distributions of several key economic parameters: claim risk, risk aversion, monitoring disutility and the three remaining utility frictions. Additional results are in the appendix: Table \ref{tab:app-est-secondary} includes additional severity and pricing latent parameters. Tables \ref{tab:app-hyperp-nonx}, \ref{tab:app-est-pricescore}, and \ref{tab:app-hyperp-x} present  all hyper-parameters outlined in Section \ref{sec:econometric-model}.

Accidents are rare: the average consumer in our panel faces  0.039  accidents per six-month period in expectation. The distribution is right-skewed:
the riskiest 5\% of drivers expect 0.113 accidents per period---about three times more than the average.  Among accidents, major ones---defined as those with Pareto losses greater than \$10{,}000---account for  10.88\%.  Combining the accident occurrence models for minor and major accidents with their respective severity models (Equations \ref{eq-claim-severity-model}), the expected cost of fully insuring a consumer is  \$205.23,  of which  \$113.62  is due to major accidents. If a consumer chooses two of the most popular plans, \$40{,}000 or \$50{,}000 coverage limits, their expected out-of-pocket expenditures would be  \$14.26  and  \$10.66, respectively.

Consumers exhibit mild risk aversion, with an average absolute risk-aversion parameter of  $1.43 \times 10^{-5}$.  At this level, an individual with average income would require a certain payment of  \$7.17  to be indifferent to a 50-50 lottery of winning or losing \$1{,}000. Across the distribution, this certainty equivalent is  \$2.09  at the 5th percentile and  \$15.09  at the 95th percentile.  Our estimate is lower than prior studies: for example, \textcite{Cohen2007} report an average of $2.6 \times 10^{-5}$, corresponding to  \$13  in the lottery interpretation.
However, this difference is expected, as our focus is on mandatory liability insurance, whereas optional coverage for one's own car tends to be purchased by more risk-averse consumers. Moreover, prior literature  has found that consumers tend to exhibit very risk averse behaviors for small risks and the opposite for large ones \parencite{sydnor2010over}. For example, \textcite{chetty2006new} estimates that, based on labor supply choices, the constant relative risk aversion parameters are around 1. At our income level, this corresponds to about \$12.12 in the lottery interpretation.

Consistent with prior literature such as \textcite{Handel2013},\footnote{The average switching cost is estimated to be \$2,032 per year.} choice frictions are economically significant. On average, consumers forgo  \$333  in potential gains per period by not switching firms, relative to an average premium of \$578. Switching plans across periods is even less common than switching firms, corresponding to an implied plan-switching cost of \$543. Similarly, the magnitude of the monitoring disutility is large: on average, consumers will only opt in if they anticipate an expected gain of at least  \$113  per period.

Although our model does not impose any explicit correlations among utility parameters, we detect meaningful ex-post correlations driven by  the correlated choice patterns in the data.  Log accident risk is strongly positively correlated with firm-switching cost ($\rho =  0.52 $) and with preference for the default minimum coverage plan ($\rho =  0.54 $).  Accident risk is weakly negatively correlated with risk aversion ($\rho =  -0.12 $) and monitoring disutility ($\rho =  -0.08 $).  At the individual level, the correlation between the private log risk component (known ex-ante to the consumer but not to the firm) with switching cost, preference for minimum coverage, risk aversion, and monitoring disutility are $\rho =  0.10,  0.02,  -0.10,  -0.22$.  Most of the correlation between risk and firm switching cost or preference for the default plan is accounted for by observable characteristics, and so their effect on selection of private risk is limited. However, privately safer drivers tend to suffer higher monitoring disutility, which may dampen advantageous selection into monitoring.

\paragraph{Robustness and Discussion}

The most binding assumption in our model is that consumers have a three-period planning horizon. In Appendix Tables \ref{tab:app-est-robust-2p} and \ref{tab:app-est-robust-4p}, we re-estimate the model under two- and four-period horizons.
Compared to the estimates in our baseline model (\Cref{tab:model_params_summary}), the main difference is that the planning horizon influences the magnitude of choice frictions needed to rationalize the observed choice patterns, beyond what can be explained by the direct certainty-equivalent value of insurance.

For instance, we estimate an average firm switching cost of \$333 under our three-period benchmark, down from \$568 with a two-period horizon and slightly above the \$329 with a four-period horizon.
The planning horizon and switching costs move in opposite directions because both affect how much today's choice ``sticks,'' and affects future-period utility.
In a three-period model, lower switching costs make the period-0 decision easier to undo at renewals, so its consequences persist less, mimicking a shorter effective horizon. When we explicitly shorten the model to two periods, we need higher switching costs to keep current choices binding and amplify their future consequences, mimicking a longer horizon.

Conversely, the one-time monitoring disutility estimate increases with the modeled planning horizon: \$113 in the baseline, versus \$90 and \$208 in the two- and four-period models, respectively.
As with switching costs, these differences in estimates make sense as the frictions are fundamentally trying to explain the same choice patterns but with different levels of insurance value to rationalize around.
Monitoring disutility is only relevant in the first period and consumers must opt-in voluntarily. A longer planning horizon typically thus implies more value from a discount and the monitoring disutility needed to rationalize not opting in is higher with a longer horizon.

These results illustrate that it may be difficult to separately identify consumers' planning horizon from the magnitude of their choice frictions. We thus do not pursue a model in which the time horizon is determined endogenously. Instead, we show that, despite differences in the magnitude of choice frictions, our main counterfactual results remain robust to alternative planning-horizon assumptions in Section \ref{subsec:model_sensitivity}.
